% ============================================================
% Problema 1
% ============================================================
\section{Problema 1: Diferentes tipos de cables}

\begin{tcolorbox}[enunciado]
Investigue sobre cuál es la diferencia entre un cable serial DTE y serial DCE.
\end{tcolorbox}

\subsection*{Solución}
Antes de mencionar la diferencia entre ambas categorías, es necesario comprender lo que significa cada una.

\begin{itemize}
    \item \textbf{¿Qué es el Data Terminal Equipment (DTE)?}
    Un Equipo Terminal de Datos es cualquier dispositivo que pueda servir como fuente o destino de datos digitales en binario. En la \textbf{capa física} puede ser una terminal, computadora, impresora o cualquier otro dispositivo que genere o consuma datos digitales. Siendo que los DTE no suelen comunicar información, pero necesitan un intermediario para poder comunicarse \cite{dte_vs_dce_cables} \cite{geeksforgeeks2020dte}.

    \item \textbf{¿Qué es el Data Circuit Terminating Equipment (DCE)?}
    Los Equipos de Comunicación de Datos son una categoría que incluye cualquier dispositivo que transmite o recibe datos en forma de señal analógica o digital a través de una red. En la \textbf{capa física}, un DCE toma datos generados por un DTE, los convierte en una señal apropiada y luego introduce la señal en un enlace de telecomunicaciones. Los DTE comúnmente utilizados en está capa incluyen módems \cite{geeksforgeeks2020dte}. 
\end{itemize}

Por lo tanto, ambas categorías de equipos (DTE y DCE) son interdependientes, es decir, que la comunicación de uno normalmente involucra al otro. Así que la \textbf{principal diferencia} entre una conexión serial \textbf{DTE} y una \textbf{DCE} radica en quién marca el ritmo (la sincronización) de la comunicación.

Como los enlaces seriales envían los datos bit por bit en una fila continua, ambos extremos necesitan estar sincronizados para saber dónde empieza y dónde termina cada bit.

El \textbf{extremo DCE} del cable es el encargado de generar y enviar la señal de reloj que dicta la velocidad a la que viajarán los datos. 
Mientas que, el \textbf{extremo DTE} del cable es quien recibe y obedece la velocidad del reloj que le manda el DCE.

\textbf{Conexiones Seriales en Cisco Packet Tracer}
\begin{figure}[H]
    \centering
    \includegraphics[width=1\linewidth]{conexiones_seriales.png}
    \caption{Captura de pantalla de la sección \textbf{Connections/Links} de la documentación del emulador Cisco Packet Tracer \cite{cisco_pt_help}.}
\end{figure}
Si se elige el tipo de conexión \textbf{Serial DCE} y, a continuación, se conectan dos dispositivos, el primer dispositivo será el lado DCE y el segundo se configurará automáticamente como el lado DTE. Lo contrario ocurre si se elige el tipo de conexión \textbf{Serial DTE}.